
\documentclass{article} % For LaTeX2e

% Optional math commands from https://github.com/goodfeli/dlbook_notation.
\input{math_commands.tex}

\usepackage[dvipsnames]{xcolor}
\usepackage{hyperref}
\hypersetup{
 colorlinks=true,
 citecolor=RoyalBlue,
 linkcolor=BrickRed,
 urlcolor=CornflowerBlue
}
\usepackage{url}
\usepackage{bbm}
\usepackage{amsmath}
\usepackage{amssymb}
\usepackage{mathtools}
\usepackage{booktabs}
\usepackage{tabularx}
\usepackage{array}
\usepackage{algorithm}
\usepackage{algpseudocode}
\usepackage{graphicx}
\usepackage{multirow}
\usepackage{colortbl,xcolor}
\usepackage{makecell}
\definecolor{indomgray}{gray}{0.9}
\newcommand{\cellgray}{\cellcolor{indomgray}}
\newcolumntype{Y}{>{\centering\arraybackslash}X}
\newcommand{\dd}{\mathrm{d}}
\def\v\epsilon{{\bm{\epsilon}}}


\title{Reinforced Flow Matching}



% The \author macro works with any number of authors. There are two commands
% used to separate the names and addresses of multiple authors: \And and \AND.
%
% Using \And between authors leaves it to \LaTeX{} to determine where to break
% the lines. Using \AND forces a linebreak at that point. So, if \LaTeX{}
% puts 3 of 4 authors names on the first line, and the last on the second
% line, try using \AND instead of \And before the third author name.

\newcommand{\fix}{\marginpar{FIX}}
\newcommand{\new}{\marginpar{NEW}}

\setlength{\parindent}{0pt}

%\iclrfinalcopy % Uncomment for camera-ready version, but NOT for submission.
\begin{document}


\maketitle

\begin{abstract}
Large language model with reinforcement learning
usually learns a reward-tilted conditional distribution,
in a parametric form.
In diffusion model,
the distribution for the clean data 
is non-parametric, and
represented by samples,
and the diffusion model learns a denoising process
to approximate the distribution.
Inspired by this,
we decompose the problem into two steps.
reweight the samples according to rewards,
and update the flow matching with weighted samples
in an incremental way.
The reweighting scheme is a non-parametric way for the reward-tilted distribution,
formulated
from KL-constrained reward maximization. 
The incremental update consists of the target velocity 
estimation from the original velocity
and velocity correction that depends on  mass shift 
and timestep-dependent contribution weight,
followed by an MSE loss for velocity update.
We also note that
several existing methods
can also be casted into 
sampling reweighting and 
velocity correction.


\end{abstract}




% \section{From rewards to sample weights}



% For a fixed prompt, let $\pi^{\mathrm{o}}(\mathbf z)$ be the old distribution
% of clean samples and $r(\mathbf z)$ their reward. We omit the prompt from the
% notation. Reward maximization with a KL penalty gives
% \begin{equation}
% \begin{gathered}
% \max_{\pi}\;
% \mathbb E_{\mathbf z\sim\pi}[r(\mathbf z)]
% -\gamma D_{\mathrm{KL}}(\pi\|\pi^{\mathrm{o}}),\\
% \pi^{\mathrm{*}}(\mathbf z)
% =\pi^{\mathrm{o}}(\mathbf z)
% \frac{\exp(r(\mathbf z)/\gamma)}
% {\mathbb E_{\mathbf z\sim\pi^{\mathrm{o}}}[\exp(r(\mathbf z)/\gamma)]},
% \end{gathered}
% \label{eq:reward-tilt}
% \end{equation}
% where $\gamma>0$ controls the strength of the reward preference. Thus, the
% distribution update can be represented by the probability ratio
% \begin{equation}
% w(\mathbf z)=\frac{\pi^{\mathrm{*}}(\mathbf z)}{\pi^{\mathrm{o}}(\mathbf z)},
% \qquad
% \mathbb E_{\mathbf z\sim\pi^{\mathrm{o}}}[w(\mathbf z)]=1.
% \label{eq:ratio}
% \end{equation}
% For $K$ rollout samples from the same prompt, we estimate these weights by
% \begin{equation}
% w(\mathbf z_k)
% =\frac{\exp(r(\mathbf z_k)/\gamma)}
% {\frac1K\sum_{j=1}^{K}\exp(r(\mathbf z_j)/\gamma)},
% \qquad
% \frac1K\sum_{k=1}^{K}w(\mathbf z_k)=1.
% \label{eq:sample-weights}
% \end{equation}
% Each sample's old empirical probability is $1/K$ and its new probability is
% $w(\mathbf z_k)/K$. The quantity $w-1$ therefore describes the relative
% increase or decrease in probability. More generally, the framework accepts
% any nonnegative weights with mean one; the exponential map is one of the choice
% which is given by Eq.~\eqref{eq:reward-tilt}.

\section{Incremental velocity target update}

\paragraph{The old velocity.}
Let $p_t(\mathbf x_t\mid\mathbf z)$ be the forward noising distribution and
$\mathbf v_t(\mathbf x_t\mid\mathbf z)$ its conditional velocity. For example,
the linear path $\mathbf x_t=(1-t)\mathbf z+t\boldsymbol\epsilon$, with
$\boldsymbol\epsilon\sim\mathcal N(\mathbf 0,\mathbf I)$, has
$\mathbf v_t(\mathbf x_t\mid\mathbf z)=(\mathbf x_t-\mathbf z)/t$ for $t>0$.
The optimal flow-matching velocity is the conditional average
\begin{align}
\mathbf v(\mathbf x_t,t)
&=\mathbb E_{\mathbf z\sim p(\cdot\mid\mathbf x_t)}
  [\mathbf v_t(\mathbf x_t\mid\mathbf z)]
\nonumber\\
&=\mathbb E_{\mathbf z\sim\pi}
  [\alpha_t(\mathbf x_t,\mathbf z)
   \mathbf v_t(\mathbf x_t\mid\mathbf z)],
\nonumber\\
\alpha_t(\mathbf x_t,\mathbf z)
&=\frac{p_t(\mathbf x_t\mid\mathbf z)}{p_t(\mathbf x_t)}.
\label{eq:old-velocity}
\end{align}
Here $p_t$ is the noisy marginal distribution, and
$\alpha_t$ measures a sample's relative contribution at
$\mathbf x_t$. For fixed $(\mathbf x_t,t)$, its expectation over
$\mathbf z\sim\pi$ is one.


Reweighting changes the clean distribution while keeping the forward noising
process fixed. Hence the new optimal velocity is
\begin{equation}
\mathbf v^{\mathrm{n}}(\mathbf x_t,t)
=\mathbb E_{z\sim\pi^{\mathrm{o}}}[w(z)\alpha_t^{\mathrm{n}}\mathbf v_t(\mathbf x_t\mid\mathbf z)].
\label{eq:new-velocity}
\end{equation}

\paragraph{One way:}
For a small incremental update, we assume that the noisy marginal distribution changes only slightly, so \textbf{we assume
$p_t^{\mathrm{n}}(\mathbf x_t)\approx p_t^{\mathrm{o}}(\mathbf x_t)$}, so
$\alpha_t^{\mathrm{n}}\approx\alpha_t^{\mathrm{o}}$. Using
$\mathbb E_{\pi^{\mathrm{o}}}[w-1]=0$ gives
\begin{align}
\mathbf v^{\mathrm{n}}(\mathbf x_t,t)
&\approx\mathbb E_{z\sim\pi^{\mathrm{o}}}[w(z)\alpha_t^{\mathrm{o}}\mathbf v_t(\mathbf x_t\mid\mathbf z)]
\nonumber\\
&=\mathbb E_{\pi^{\mathrm{o}}}
   [\alpha_t^{\mathrm{o}}\mathbf v_t(\mathbf x_t\mid\mathbf z)]
 +\mathbb E_{\pi^{\mathrm{o}}}[(w-1)\alpha_t^{\mathrm{o}}\mathbf v_t(\mathbf x_t\mid\mathbf z)]\\
&=\mathbf v^{\mathrm{o}}(\mathbf x_t,t)
 +\mathbb E_{\pi^{\mathrm{o}}}[(w-1)\alpha_t^{\mathrm{o}}\mathbf v_t(\mathbf x_t\mid\mathbf z)]
\nonumber\\
&=\mathbf v^{\mathrm{o}}(\mathbf x_t,t)
 +\mathbb E_{\pi^{\mathrm{o}}}
   [(w-1)(\alpha_t^{\mathrm{o}}\mathbf v_t(\mathbf x_t\mid\mathbf z)-\mathbf v^{\mathrm{o}}(\mathbf x_t,t))].
\label{eq:increment}
\end{align}
% Replacing this expectation with average gives
% \begin{equation}
% \boxed{
% \mathbf v^{\mathrm{n}}(\mathbf x_t,t)
% \approx\mathbf v^{\mathrm{o}}(\mathbf x_t,t)
% +\frac1K\sum_{k=1}^{K}(w(\mathbf z_k)-1)
% \alpha_t^{\mathrm{o}}(\mathbf x_t,\mathbf z_k)
% \left(
% \mathbf v_t(\mathbf x_t\mid\mathbf z_k)
% -\frac{\mathbf v^{\mathrm{o}}(\mathbf x_t,t)}
% {\alpha_t^{\mathrm{o}}(\mathbf x_t,\mathbf z_k)}
% \right).
% }
% \label{eq:rollout-increment}
% \end{equation}

\paragraph{Another way:}

Since $w=\pi^{\mathrm{n}}/\pi^{\mathrm{o}}$,
\begin{align}
\mathbb E_{\mathbf z\sim\pi^{\mathrm{o}}}[w\alpha_t^{\mathrm{o}}]
&=\mathbb E_{\mathbf z\sim\pi^{\mathrm{o}}}[\frac{\pi^{\mathrm{n}}}{\pi^{\mathrm{o}}}\frac{p_t(\mathbf x_t\mid\mathbf z)}{p_t^{\mathrm{o}}(\mathbf x_t)}]\\
&=\mathbb E_{\mathbf z\sim\pi^{\mathrm{n}}}[\frac{p_t(\mathbf x_t\mid\mathbf z)}{p_t^{\mathrm{o}}(\mathbf x_t)}]\\
&=\frac{p_t^{\mathrm{n}}(\mathbf x_t)}{p_t^{\mathrm{o}}(\mathbf x_t)},
\label{eq:velocity-normalizer}
\end{align}
So, 
\begin{align}
\mathbf v^{\mathrm{n}}(\mathbf x_t,t)
&=\mathbb E_{\mathbf z\sim\pi^{\mathrm{o}}}
 [w\alpha_t^{\mathrm{n}}\mathbf v_t(\mathbf x_t\mid\mathbf z)]
\nonumber\\
&=\frac{\mathbb E_{\mathbf z\sim\pi^{\mathrm{o}}}
 [w\alpha_t^{\mathrm{o}}\mathbf v_t(\mathbf x_t\mid\mathbf z)]}
 {\mathbb E_{\mathbf z\sim\pi^{\mathrm{o}}}[w\alpha_t^{\mathrm{o}}]}.
\label{eq:normalized-new-velocity}
\end{align}
The old velocity satisfies
$\mathbb E_{\mathbf z\sim\pi^{\mathrm{o}}}[\alpha_t^{\mathrm{o}}\mathbf v_t(\mathbf x_t\mid\mathbf z)]
=\mathbf v^{\mathrm{o}}(\mathbf x_t,t)$ and
$\mathbb E_{\mathbf z\sim\pi^{\mathrm{o}}}[\alpha_t^{\mathrm{o}}]=1$, so
\begin{equation}
\mathbb E_{\mathbf z\sim\pi^{\mathrm{o}}}
 [\alpha_t^{\mathrm{o}}(\mathbf v_t(\mathbf x_t\mid\mathbf z)-\mathbf v^{\mathrm{o}}(\mathbf x_t,t))]=0.
\label{eq:centered-old-velocity}
\end{equation}
Subtracting this zero term from the numerator gives the exact increment:
\begin{align}
&\mathbf v^{\mathrm{n}}(\mathbf x_t,t)-\mathbf v^{\mathrm{o}}(\mathbf x_t,t)
\nonumber\\
&\quad=\frac{\mathbb E_{\mathbf z\sim\pi^{\mathrm{o}}}
 [w\alpha_t^{\mathrm{o}}\mathbf v_t(\mathbf x_t\mid\mathbf z)]}
 {\mathbb E_{\mathbf z\sim\pi^{\mathrm{o}}}[w\alpha_t^{\mathrm{o}}]} 
-\mathbf v^{\mathrm{o}}(\mathbf x_t,t)
\nonumber\\
&\quad=\frac{\mathbb E_{\mathbf z\sim\pi^{\mathrm{o}}}
 [w\alpha_t^{\mathrm{o}}(\mathbf v_t(\mathbf x_t\mid\mathbf z)-\mathbf v^{\mathrm{o}}(\mathbf x_t,t))]}
 {\mathbb E_{\mathbf z\sim\pi^{\mathrm{o}}}[w\alpha_t^{\mathrm{o}}]}
\nonumber\\
&\quad=\frac{\mathbb E_{\mathbf z\sim\pi^{\mathrm{o}}}
 [(w-1)\alpha_t^{\mathrm{o}}(\mathbf v_t(\mathbf x_t\mid\mathbf z)-\mathbf v^{\mathrm{o}}(\mathbf x_t,t))]}
 {\mathbb E_{\mathbf z\sim\pi^{\mathrm{o}}}[w\alpha_t^{\mathrm{o}}]}
\nonumber\\
&\quad=\frac{p_t^{\mathrm{o}}(\mathbf x_t)}{p_t^{\mathrm{n}}(\mathbf x_t)}
 \mathbb E_{\mathbf z\sim\pi^{\mathrm{o}}}
 [(w-1)\alpha_t^{\mathrm{o}}(\mathbf v_t(\mathbf x_t\mid\mathbf z)-\mathbf v^{\mathrm{o}}(\mathbf x_t,t))].
\label{eq:exact-residual-increment}
\end{align}
With the same small-change assumption,
\begin{equation}
\frac{p_t^{\mathrm{n}}(\mathbf x_t)}{p_t^{\mathrm{o}}(\mathbf x_t)}
\approx1,
\label{eq:small-marginal-change}
\end{equation}
We have
\begin{equation}
\begin{aligned}
\mathbf v^{\mathrm{n}}(\mathbf x_t,t)
\approx\mathbf v^{\mathrm{o}}(\mathbf x_t,t)
+\mathbb E_{\mathbf z\sim\pi^{\mathrm{o}}}
 [(w-1)\alpha_t^{\mathrm{o}}(\mathbf v_t(\mathbf x_t\mid\mathbf z)-\mathbf v^{\mathrm{o}}(\mathbf x_t,t))].
\end{aligned}
\label{eq:residual-increment}
\end{equation}
This separates three roles: $w-1$ sets the probability change,
$\alpha_t^{\mathrm{o}}$ measures the contribution at the current noisy state,
and $\mathbf v_t(\mathbf x_t\mid\mathbf z)-\mathbf v^{\mathrm{o}}(\mathbf x_t,t)$ gives the correction direction.
A weight above one pulls toward that sample's velocity; a weight below one
pushes away. A weight of one contributes no correction.


% And we also replace this expectation with average
% \begin{equation}
% \boxed{
% \mathbf v^{\mathrm{n}}(\mathbf x_t,t)
% \approx\mathbf v^{\mathrm{o}}(\mathbf x_t,t)
% +\frac1K\sum_{k=1}^{K}(w(\mathbf z_k)-1)
% \alpha_t^{\mathrm{o}}(\mathbf x_t,\mathbf z_k)
% \left(
% \mathbf v_t(\mathbf x_t\mid\mathbf z_k)
% -\mathbf v^{\mathrm{o}}(\mathbf x_t,t)
% \right).
% }
% \label{eq:rollout-residual-increment}
% \end{equation}

\paragraph{Estimating the contribution.}
The marginal density in $\alpha_t$ is expensive to compute. We therefore
use a normalized velocity error as a estimate:
\begin{equation}
\begin{aligned}
\widetilde\alpha_t(\mathbf x_t,\mathbf z)
&=\frac{1}{\operatorname{mean}
 |\mathbf v_t(\mathbf x_t\mid\mathbf z)-\mathbf v_\theta(\mathbf x_t,t)|},\\
\alpha_t(\mathbf x_t,\mathbf z)
&\approx\frac{\widetilde\alpha_t(\mathbf x_t,\mathbf z)}
{\frac1B\sum_{b=1}^{B}
 \widetilde\alpha_{t_b}(\mathbf x_{t_b},\mathbf z_b)}.
\end{aligned}
\label{eq:alpha-estimate}
\end{equation}
The mean in the error is over feature dimensions, $B$ is the batch size. This assigns larger
contributions to samples whose conditional velocities are closer to the
model prediction.

\section{Gap between current implemented algorithm and our derivation}

With gradient accumulation over all mini batches before one optimizer
update, shuffling samples only changes how the computation is split.
The accumulated gradient corresponds to the
mean loss over all $CK$ samples. For simplicity, we consider a single prompt and its
$K$ samples group below.



\paragraph{Equivalent v-prediction policy loss of the current implementation.}
Let $\mathbf v_i=\mathbf v_{t_i}(\mathbf x_i^{t_i}\mid\mathbf x_i^0)$
be the conditional velocity, and let $\mathbf v_i^\theta$ and
$\mathbf v_i^{\mathrm{o}}$ be the current and old predictions.
Using $\mathbf x_i^\theta=\mathbf x_i^{t_i}-t_i\mathbf v_i^\theta$
and the same relation for the old prediction, the algorithm's policy loss
is exactly
\begin{align}
\mathbf v_i^{\mathrm{target}}
&=\mathbf v_i^{\mathrm{o}}
+\beta(w_i-1)(\mathbf v_i-\mathbf v_i^{\mathrm{o}}),
\label{eq:implemented-v-target}\\
\mathcal L_{\mathrm{policy}}^{\mathrm{impl}}
&=\frac1K\sum_{i=1}^{K}
\frac{t_i\|\mathbf v_i^\theta
-\mathbf v_i^{\mathrm{target}}\|_2^2}
{\operatorname{sg}[\operatorname{mean}
|\mathbf v_i-\mathbf v_i^{\mathrm{o}}|]}.
\label{eq:implemented-v-loss}
\end{align}


\paragraph{Policy loss with the derived target.}
Using $\mathbf v^{\mathrm{n}}$ from Eq.~\eqref{eq:residual-increment}
as the target and keeping a timestep weight $t_i$ outside MSE, we define
\begin{equation}
\boxed{
\begin{aligned}
\mathbf v_i^{\mathrm{target}}
&=\mathbf v^{\mathrm{o}}_i
+\mathbb E_{\mathbf z\sim\pi^{\mathrm{o}}}
 [(w_i-1)\alpha_{t_i}^{\mathrm{o}}(\mathbf v_i-\mathbf v^{\mathrm{o}}_i)],\\
\mathcal L_{\mathrm{policy}}^{\mathrm{derived}}
&=\frac1K\sum_{i=1}^{K}t_i
\left\|\mathbf v_i^\theta
-\mathbf v_i^{\mathrm{target}}\right\|_2^2.
\end{aligned}
}
\label{eq:derived-policy-loss}
\end{equation}




\begin{algorithm}[H]
    \caption{Single Target Loss Base Scheme }
    \label{alg:velocity-add-delete}
    \textbf{Require:} Initial model $\theta_0$, reward function $R$, prompt distribution $p(c)$, number of prompts per iteration $C$, rollout size per prompt $K$, batch size $B$, total samples per iteration $D = KC$, rollout coverage coefficient $\beta$, mass-shift clip $\Delta_{\mathrm{clip}}$, reference regularization $\lambda_{\mathrm{ref}}$, EMA rate $\eta$, reward-scale constant $\epsilon_R$, numerical constant $\epsilon_{\mathrm{alg}}$, gradient accumulation steps $N_{\mathrm{acc}}$.\\
    \textbf{Initialize:} Training policy $\theta \leftarrow \theta_0$, old policy $\theta_{\mathrm{old}} \leftarrow \theta_0$, reference policy $\theta_{\mathrm{ref}} \leftarrow \theta_0$, data buffer $\mathcal{D} \leftarrow \emptyset$.
    \begin{algorithmic}[1]
        \For{\text{each iteration}}
            \For{\text{each prompt $c_j \sim p(c)$, $j = 1, \dots, C$}} \Comment{Rollout Step, Data Collection}
                \State Sample $K$ clean images $\{\mathbf{x}_0^{j,k}\}_{k=1}^{K}$ with $\theta$ conditioned on $c_j$.
                \State Evaluate rewards $R_{j,k} \leftarrow R(\mathbf{x}_0^{j,k}, c_j)$ for $k = 1, \dots, K$.
            \EndFor

            \State Compute prompt-wise reward means
            $\bar{R}_{c_j} \leftarrow \frac{1}{K}\sum_{k=1}^{K} R_{j,k}$ for $j = 1, \dots, C$.

            \State Compute global reward mean
            $\bar{R} \leftarrow \frac{1}{D}\sum_{j=1}^{C}\sum_{k=1}^{K} R_{j,k}$.

            \State Compute global reward scale
            $\sigma_R \leftarrow \sqrt{\frac{1}{D}\sum_{j=1}^{C}\sum_{k=1}^{K}(R_{j,k} - \bar{R})^2}.$

            \For{\text{each prompt $c_j$, $j = 1, \dots, C$}}
                \State Compute normalized weights
                \[
                    \Delta_{j,k} \leftarrow \frac{1}{\Delta_{\mathrm{clip}}}
                    \operatorname{clip}\!\left(
                        \frac{R_{j,k} - \bar{R}_{c_j}}{\sigma_R},
                        -\Delta_{\mathrm{clip}},
                        \Delta_{\mathrm{clip}}
                    \right),
                    \quad
                    w_{j,k}-1\leftarrow \Delta_{j,k},
                    \quad k = 1, \dots, K.
                \]

                \State Store
                $\mathcal{D}\leftarrow\mathcal{D}\cup
                \{(c_j,\mathbf{x}_0^{j,k},w_{j,k})\}_{k=1}^{K}$.
            \EndFor

            \For{\text{each mini batch }
            $\{(c_i,\mathbf{x}_i^0,w_i)\}_{i=1}^{B}\subset\mathcal{D}$}
                \State Sample $t_i\sim\mathcal{U}(0,1)$,
                $\boldsymbol{\epsilon}_i\sim\mathcal{N}(\mathbf{0},\mathbf{I})$;
                set $\mathbf{x}_i^{t_i}\leftarrow
                (1-t_i)\mathbf{x}_i^0+t_i\boldsymbol{\epsilon}_i$.
                \State Predict $\mathbf{x}_i^\theta$,
                $\mathbf{x}_i^{\mathrm{old}}$, and $\mathbf{x}_i^{\mathrm{ref}}$
                from $(\mathbf{x}_i^{t_i},c_i,t_i)$ with
                $\theta$, $\theta_{\mathrm{old}}$, and $\theta_{\mathrm{ref}}$.
                \State $\mathbf{x}_i^{\mathrm{target}}\leftarrow
                \mathbf{x}_i^{\mathrm{old}}+
                \beta(w_i-1)(\mathbf{x}_i^0-\mathbf{x}_i^{\mathrm{old}})$.
                \State Compute
                \[
                    \begin{aligned}
                    \mathcal{L}_{\mathrm{policy}}
                    &\leftarrow\frac{1}{B}\sum_{i=1}^{B}
                    \frac{\|\mathbf{x}_i^\theta-\mathbf{x}_i^{\mathrm{target}}\|_2^2}
                    {\operatorname{sg}[\operatorname{mean}
                    |\mathbf{x}_i^0-\mathbf{x}_i^{\mathrm{old}}|]},\\
                    \mathcal{L}_{\mathrm{ref}}
                    &\leftarrow\frac{1}{B}\sum_{i=1}^{B}
                    \|\mathbf{v}_i^\theta-\mathbf{v}_i^{\mathrm{ref}}\|_2^2.
                    \end{aligned}
                \]
                \State Backpropagate
                $(\Delta_{clip}\mathcal{L}_{\mathrm{policy}}+
                \lambda_{\mathrm{ref}}\mathcal{L}_{\mathrm{ref}})/N_{\mathrm{acc}}$.
            \EndFor
            \State Update $\theta$; set
            $\theta_{\mathrm{old}}\leftarrow
            \eta\theta_{\mathrm{old}}+(1-\eta)\theta$ and
            $\mathcal{D}\leftarrow\emptyset$.
        \EndFor
    \end{algorithmic}
    \textbf{Output:} $\mathbf{v}_\theta$
\end{algorithm}







\section{Reweighting design}
\label{sec:group-reweighting}


\begin{algorithm}[H]
    \scriptsize
    \setlength{\abovedisplayskip}{1pt}
    \setlength{\belowdisplayskip}{1pt}
    \renewcommand{\arraystretch}{0.9}
    \caption{Selection scheme: Asymmetric Selection}
    \label{alg:x-prediction-asymmetric-hard-selection}
    \textbf{Require:} Initial model $\theta_0$, reward function $R$,
    prompt distribution $p(c)$, $C$ prompts per iteration, $K$ samples per
    prompt, $D=CK$, batch size $B$, retained ratio $\rho$, shift strength
    $\beta$, mass-shift clip $\Delta_{\mathrm{clip}}$, reference coefficient
    $\lambda_{\mathrm{ref}}$, EMA rate $\eta$,
    accumulation steps $N_{\mathrm{acc}}$.\\
    \textbf{Initialize:}
    $\theta,\theta_{\mathrm{old}},\theta_{\mathrm{ref}}\leftarrow\theta_0$,
    $\mathcal{D}\leftarrow\emptyset$.
    \begin{algorithmic}[1]
        \For{\text{each iteration}}
            \For{\text{each prompt $c_j\sim p(c)$, $j=1,\dots,C$}}
                \State Sample $\{\mathbf{x}_0^{j,k}\}_{k=1}^{K}$ with $\theta$
                conditioned on $c_j$ and set
                $R_{j,k}\leftarrow R(\mathbf{x}_0^{j,k},c_j)$.

            \EndFor

            \State Compute global reward mean
            $\bar{R}\leftarrow \frac{1}{D}\sum_{j=1}^{C}\sum_{k=1}^{K}R_{j,k}$.
            \State Compute global reward scale
            $\sigma_R \leftarrow \sqrt{\frac{1}{D}\sum_{j=1}^{C}\sum_{k=1}^{K}(R_{j,k} - \bar{R})^2}.$

            \For{\text{each prompt $c_j\sim p(c)$, $j=1,\dots,C$}}

                \State Set the mass-shift center to the prompt-group mean:
                $b_j\leftarrow K^{-1}\sum_{k=1}^{K}R_{j,k}$.
                \State Within this prompt group, compute
                \[
                    % a_j^-\leftarrow\max_{R_{j,k}<b_j}|R_{j,k}-b_j|,
                    % \quad
                    % a_j^+\leftarrow\max_{R_{j,k}>b_j}|R_{j,k}-b_j|,
                    % \quad
                    \Delta_{j,k}\leftarrow
                    \frac{1}{\Delta_{\mathrm{clip}}}
                    \operatorname{clip}\!\left(
                        \frac{R_{j,k}-b_j}{\sigma_R},
                        -\Delta_{\mathrm{clip}},
                        \Delta_{\mathrm{clip}}
                    \right)
                \]
                \State Sort the $K$ samples so that
                $\Delta_{j,1}\le\cdots\le\Delta_{j,q_j}\le0
                \le\Delta_{j,q_j+1}\le\cdots\le\Delta_{j,K}$.
                \State Select the positive strong-tail boundary:
                \[
                    M_j^+\leftarrow\sum_{k=q_j+1}^{K}\Delta_{j,k},
                    \qquad
                    h_{j,\rho}\leftarrow
                    \arg\min_{q_j+1\le h\le K}
                    \left|\sum_{k=h}^{K}\Delta_{j,k}-\rho M_j^+\right|.
                \]
                \State Keep all negative samples and the positive strong signal samples:
                \[
                    s_{j,k}\leftarrow
                    \begin{cases}
                        \Delta_{j,k},&k\le q_j,\\
                        0,&q_j<k<h_{j,\rho},\\
                        \Delta_{j,k},&k\ge h_{j,\rho}.
                    \end{cases}
                \]
                \State Normalize the selected mass within this prompt group:
                \[
                    % m_j^-\leftarrow\sum_{s_{j,k}<0}(-s_{j,k}),\quad
                    % m_j^+\leftarrow\sum_{s_{j,k}>0}s_{j,k},\quad
                    % w_{j,k}-1\leftarrow
                    % \begin{cases}
                    %     s_{j,k}/m_j^-,&s_{j,k}<0,\\
                    %     0,&s_{j,k}=0,\\
                    %     s_{j,k}/m_j^+,&s_{j,k}>0.
                    % \end{cases}
                    w_{j,k}-1\leftarrow
                    \begin{cases}
                        s_{j,k},&s_{j,k}<0,\\
                        0,&s_{j,k}=0,\\
                        s_{j,k}/\rho,&s_{j,k}>0.
                    \end{cases}                    
                \]
                \State Store
                $\mathcal{D}\leftarrow\mathcal{D}\cup
                \{(c_j,\mathbf{x}_0^{j,k},w_{j,k})\}_{k=1}^{K}$.
            \EndFor

            \For{\text{each mini batch }
            $\{(c_i,\mathbf{x}_i^0,w_i)\}_{i=1}^{B}\subset\mathcal{D}$}
                \State Sample $t_i\sim\mathcal{U}(0,1)$,
                $\boldsymbol{\epsilon}_i\sim\mathcal{N}(\mathbf{0},\mathbf{I})$;
                set $\mathbf{x}_i^{t_i}\leftarrow
                (1-t_i)\mathbf{x}_i^0+t_i\boldsymbol{\epsilon}_i$.
                \State Predict $\mathbf{x}_i^\theta$,
                $\mathbf{x}_i^{\mathrm{old}}$, and $\mathbf{x}_i^{\mathrm{ref}}$
                from $(\mathbf{x}_i^{t_i},c_i,t_i)$ with
                $\theta$, $\theta_{\mathrm{old}}$, and $\theta_{\mathrm{ref}}$.
                \State $\mathbf{x}_i^{\mathrm{target}}\leftarrow
                \mathbf{x}_i^{\mathrm{old}}+
                \beta(w_i-1)(\mathbf{x}_i^0-\mathbf{x}_i^{\mathrm{old}})$.
                \State Compute
                \[
                    \begin{aligned}
                    \mathcal{L}_{\mathrm{policy}}
                    &\leftarrow\frac{1}{B}\sum_{i=1}^{B}
                    \frac{\|\mathbf{x}_i^\theta-\mathbf{x}_i^{\mathrm{target}}\|_2^2}
                    {\operatorname{sg}[\operatorname{mean}
                    |\mathbf{x}_i^0-\mathbf{x}_i^{\mathrm{old}}|]},\\
                    \mathcal{L}_{\mathrm{ref}}
                    &\leftarrow\frac{1}{B}\sum_{i=1}^{B}
                    \|\mathbf{v}_i^\theta-\mathbf{v}_i^{\mathrm{ref}}\|_2^2.
                    \end{aligned}
                \]
                \State Backpropagate
                $(\Delta_{\mathrm{clip}}\mathcal{L}_{\mathrm{policy}}+
                \lambda_{\mathrm{ref}}\mathcal{L}_{\mathrm{ref}})/N_{\mathrm{acc}}$.
            \EndFor
            \State Update $\theta$; set
            $\theta_{\mathrm{old}}\leftarrow
            \eta\theta_{\mathrm{old}}+(1-\eta)\theta$ and
            $\mathcal{D}\leftarrow\emptyset$.
        \EndFor
    \end{algorithmic}
    \textbf{Output:} $\mathbf{x}_\theta$
\end{algorithm}


We now use Algorithm~\ref{alg:x-prediction-asymmetric-hard-selection} to study how raw rewards
should determine the probability ratio $w$. For one prompt, let $R_1,\ldots,R_K$ be the rewards
of its $K$ samples. Consider the group-wise empirical distribution, the old probability of
each sample is $1/K$. A nonnegative weight with group mean equal to one defines
the new probability $w_i/K$ for each sample, so
\begin{equation}
w_i=\frac{\pi^{\mathrm{n}}(\mathbf z_i)}
{\pi^{\mathrm{o}}(\mathbf z_i)},
\qquad w_i\ge0,\qquad \frac1K\sum_{i=1}^{K}w_i=1.
\label{eq:group-probability-ratio}
\end{equation}
The current target uses $\beta(w_i-1)$ in
Eq.~\eqref{eq:implemented-v-target}, and $\beta$ control the step size of target shift. 

\subsection{Concentrate the confident positive updates}

\paragraph{Base signal.}
Let $\bar R=K^{-1}\sum_{i=1}^{K}R_i$:
\begin{equation}
\Delta_i=\frac{1}{\Delta_{\mathrm{clip}}}
\operatorname{clip}\!\left(
\frac{R_i-\bar R}{\sigma_R},
-\Delta_{\mathrm{clip}},\Delta_{\mathrm{clip}}\right).
\label{eq:selection-base-signal}
\end{equation}
Here $\bar R$ is the prompt-group mean, while $\sigma_R$ is the reward
standard deviation over the full rollout samples, as in the DiffusionNFT
implementation. The base rule is $w_i-1=\Delta_i$. Positive signals pull
toward above-average samples and negative signals push away from
below-average samples linearly.

\paragraph{Selection scheme: Keep all negative signals, select strong positive signals.}
The idea is to move positive update strength from slightly above-average
samples to better samples. Sort the positive $\Delta_i$ from largest to
smallest and retain a fraction $\rho\in(0,1]$ of their total signal mass.
Let $s_i$ be the retained positive signal. We keep
samples up to the boundary closest to the desired mass. Stronger positive samples have $s_i=\Delta_i$; weaker ones have $s_i=0$. So that
\begin{equation}
\sum_{i:\Delta_i>0}s_i
=\rho\sum_{i:\Delta_i>0}\Delta_i.
\label{eq:selection-retained-mass}
\end{equation}

Finally sets
\begin{equation}
w_i-1=
\begin{cases}
\Delta_i, & \Delta_i\le0,\\
s_i/\rho, & \Delta_i>0.
\end{cases}
\label{eq:scheme-c-update}
\end{equation}
Dividing by $\rho$ restores the total positive signal after
selection. The negative side is unchanged. When $\rho=1$, this reduces to
the base scheme. An unselected positive sample has $w_i=1$: its target
stays at the old prediction. This scheme tests whether
positive updates benefit from more selective allocation.

\paragraph{Evidence at larger $\beta$.}


\begin{figure}[H]
    \centering
    \includegraphics[width=1\linewidth]{imgs/ablation-selection-setting-search-for-schemeC.png}
    \caption{Ablation for selection scheme on Pickscore.}
    \label{fig:placeholder}
\end{figure}

\newpage





\subsection{From simple selection to reward mappings}




\begin{algorithm}[H]
    \caption{Exploration of reward mappings to probability ratio $w$}
    \label{alg:velocity-add-delete}
    \textbf{Require:} Initial model $\theta_0$, reward function $R$, prompt distribution $p(c)$, number of prompts per iteration $C$, rollout size per prompt $K$, batch size $B$, total samples per iteration $D = KC$, rollout coverage coefficient $\beta$, mass-shift clip $\Delta_{\mathrm{clip}}$, reference regularization $\lambda_{\mathrm{ref}}$, EMA rate $\eta$, reward-scale constant $\epsilon_R$, numerical constant $\epsilon_{\mathrm{alg}}$, gradient accumulation steps $N_{\mathrm{acc}}$.\\
    \textbf{Initialize:} Training policy $\theta \leftarrow \theta_0$, old policy $\theta_{\mathrm{old}} \leftarrow \theta_0$, reference policy $\theta_{\mathrm{ref}} \leftarrow \theta_0$, data buffer $\mathcal{D} \leftarrow \emptyset$.
    \begin{algorithmic}[1]
        \For{\text{each iteration}}
            \For{\text{each prompt $c_j \sim p(c)$, $j = 1, \dots, C$}} \Comment{Rollout Step, Data Collection}
                \State Sample $K$ clean images $\{\mathbf{x}_0^{j,k}\}_{k=1}^{K}$ with $\theta$ conditioned on $c_j$.
                \State Evaluate rewards $R_{j,k} \leftarrow R(\mathbf{x}_0^{j,k}, c_j)$ for $k = 1, \dots, K$.
            \EndFor

            % \State Compute prompt-wise reward means
            % $\bar{R}_{c_j} \leftarrow \frac{1}{K}\sum_{k=1}^{K} R_{j,k}$ for $j = 1, \dots, C$.

            % \State Compute global reward mean
            % $\bar{R} \leftarrow \frac{1}{D}\sum_{j=1}^{C}\sum_{k=1}^{K} R_{j,k}$.

            % \State Compute global reward scale
            % $\sigma_R \leftarrow \sqrt{\frac{1}{D}\sum_{j=1}^{C}\sum_{k=1}^{K}(R_{j,k} - \bar{R})^2}.$

            \For{\text{each prompt $c_j$, $j = 1, \dots, C$}}
                \State Compute normalized weights
                \[
                    w_{j,k} \leftarrow \frac{f(R_{j,k})}{\frac{1}{K}\sum_{k=1}^{K}f(R_{j,k})},
                    \quad k = 1, \dots, K.
                \]

                \State Store
                $\mathcal{D}\leftarrow\mathcal{D}\cup
                \{(c_j,\mathbf{x}_0^{j,k},w_{j,k})\}_{k=1}^{K}$.
            \EndFor

            \For{\text{each mini batch }
            $\{(c_i,\mathbf{x}_i^0,w_i)\}_{i=1}^{B}\subset\mathcal{D}$}
                \State Sample $t_i\sim\mathcal{U}(0,1)$,
                $\boldsymbol{\epsilon}_i\sim\mathcal{N}(\mathbf{0},\mathbf{I})$;
                set $\mathbf{x}_i^{t_i}\leftarrow
                (1-t_i)\mathbf{x}_i^0+t_i\boldsymbol{\epsilon}_i$.
                \State Predict $\mathbf{x}_i^\theta$,
                $\mathbf{x}_i^{\mathrm{old}}$, and $\mathbf{x}_i^{\mathrm{ref}}$
                from $(\mathbf{x}_i^{t_i},c_i,t_i)$ with
                $\theta$, $\theta_{\mathrm{old}}$, and $\theta_{\mathrm{ref}}$.
                \State $\mathbf{x}_i^{\mathrm{target}}\leftarrow
                \mathbf{x}_i^{\mathrm{old}}+
                \beta(w_i-1)(\mathbf{x}_i^0-\mathbf{x}_i^{\mathrm{old}})$.
                \State Compute
                \[
                    \begin{aligned}
                    \mathcal{L}_{\mathrm{policy}}
                    &\leftarrow\frac{1}{B}\sum_{i=1}^{B}
                    \frac{\|\mathbf{x}_i^\theta-\mathbf{x}_i^{\mathrm{target}}\|_2^2}
                    {\operatorname{sg}[\operatorname{mean}
                    |\mathbf{x}_i^0-\mathbf{x}_i^{\mathrm{old}}|]},\\
                    \mathcal{L}_{\mathrm{ref}}
                    &\leftarrow\frac{1}{B}\sum_{i=1}^{B}
                    \|\mathbf{v}_i^\theta-\mathbf{v}_i^{\mathrm{ref}}\|_2^2.
                    \end{aligned}
                \]
                \State Backpropagate
                $(\Delta_{clip}\mathcal{L}_{\mathrm{policy}}+
                \lambda_{\mathrm{ref}}\mathcal{L}_{\mathrm{ref}})/N_{\mathrm{acc}}$.
            \EndFor
            \State Update $\theta$; set
            $\theta_{\mathrm{old}}\leftarrow
            \eta\theta_{\mathrm{old}}+(1-\eta)\theta$ and
            $\mathcal{D}\leftarrow\emptyset$.
        \EndFor
    \end{algorithmic}
    \textbf{Output:} $\mathbf{v}_\theta$
\end{algorithm}



We extend the selection idea by using a reward map $f$ to define the new
probability of each sample. For a fixed prompt $c_j$, consider its $K$
samples with rewards $R_{j,1},\ldots,R_{j,K}$. As in
Eq.~\eqref{eq:group-probability-ratio}, we work with the empirical group
distribution and omit the prompt conditioning from $\pi$:
\begin{equation}
\pi^{\mathrm{o}}(\mathbf z_i)=\frac1K,
\qquad
\pi^{\mathrm{n}}(\mathbf z_i)
=\frac{f(R_{j,i})}{\sum_{\ell=1}^{K}f(R_{j,\ell})}.
\end{equation}
The corresponding probability ratio is therefore
\begin{equation}
w_{j,i}
=\frac{\pi^{\mathrm{n}}(\mathbf z_i)}{\pi^{\mathrm{o}}(\mathbf z_i)}
=\frac{f(R_{j,i})}{\frac1K\sum_{\ell=1}^{K}f(R_{j,\ell})}.
\label{eq:normalized-reward-map}
\end{equation}
For nonnegative $f$ with a positive sum, this construction gives
$w_{j,i}\ge0$ and $K^{-1}\sum_i w_{j,i}=1$. Thus, $w_{j,i}$ is a
probability ratio, while $w_{j,i}/K$ is the new sample probability.
The signed change $w_{j,i}-1$ controls the target correction:
$w_{j,i}>1$ increases the sample's probability in the desired distribution,
$w_{j,i}<1$ decreases it, and $w_{j,i}=1$ leaves it unchanged.
If all mapped rewards are zero, we keep the old uniform distribution
by setting $w_{j,i}=1$ for all samples.

For mappings that depend on the reward statistics of each prompt group,
we define
\[
\bar R_{c_j}=\frac1K\sum_{i=1}^{K}R_{j,i},
\qquad
\sigma_{c_j}=\sqrt{\frac1K\sum_{i=1}^{K}
\bigl(R_{j,i}-\bar R_{c_j}\bigr)^2}.
\]
Unless stated otherwise, the experiments below normalize mapped rewards
within each prompt using
Eq.~\eqref{eq:normalized-reward-map}.



\paragraph{Motivation and experiment plan.}
We want to understand why the softmax weights in FlowAWR can lead to fast
convergence. Under \eqref{eq:normalized-reward-map}, they give more of each group's total
weight to samples with higher rewards. We test whether this focus on better
samples is the main reason.

Experiment A pushes this idea to the extreme situation. A one-hot rule gives all weight to
the top sample(s) and zero to the rest. This sharp choice may make training
unstable, so we also give weight to samples within \(cR_j^{\max}\) of the top
reward, without using the group standard deviation for this threshold.
If no reward in the group is positive, we use uniform weights and apply no
reward-driven target shift to that group.

Experiment B straightly tests the idea with softmax. Compared with \(c=1\),
\(c<1\) puts more weight on high-reward samples, while \(c>1\) spreads weight
more evenly. If \(c<1\) works better, it supports our idea.

Experiment C uses the unshifted linear mapping as a baseline. Experiment D
compares sigmoid with linear within each group: depending on \(c\) and the
rewards, sigmoid can yield either sharper or flatter normalized weights.
We test whether smaller \(c\) helps, whether \(c=0.3\) underperforms, and
whether a small-\(c\) sigmoid is preferable to the softmax-style mappings.

Experiment E changes the power \(p\) to
test whether steeper or flatter maps work better. It is another view to test our idea.

Experiment F uses an L1
scale. Its weight curve looks much like the softmax curve in Exp.B, so it is
less important and may be skipped.



\paragraph{Experiment A: One-hot and near-maximum threshold.}

Let
\[
R_j^{\max}=\max_i R_{j,i},\qquad
T_j(c)=R_j^{\max}-cR_j^{\max}=(1-c)R_j^{\max}.
\]
\[
S_j(c)=\{i:R_{j,i}\ge T_j(c)\}\qquad\text{for }R_j^{\max}>0.
\]
Set
\[
w_{j,i}(c)
\;=\;
\begin{cases}
\displaystyle\frac{K}{|S_j(c)|}\,\mathbbm{1}[i\in S_j(c)],
& R_j^{\max}>0,\\[6pt]
1, & R_j^{\max}\le0,
\end{cases}
\qquad c\in\{0,0.05,0.1,0.25\}.
\]
For \(R_j^{\max}>0\), when the maximum is unique, \(c=0\) is one-hot with weight \(K\) on the
max-reward sample. If several samples are equal-reward at the maximum, they share the
K mass equally. Increasing \(c\) includes samples
within a larger fraction of the maximum reward. 

For \(R_j^{\max}\le0\),
all $w_{j,i}(c)$ equal \(1\), so the probability for each sample in this group $c_j$ will not change. While other loss terms like mse between $v_{\theta}$ and $v_{old}$ and KL terms still apply.

\textbf{Target}: Test whether assigning most or all of each group's weight
mass to its best-reward samples is the main source of the gain from
softmax weighting.
% \subparagraph{Soft Top-$k$ with boundary ties.}
% As a complementary fixed-size concentration rule, solve
% \[
% \max_{\{w_{j,i}\}}\sum_{i=1}^{K}w_{j,i}R_{j,i}
% \quad\text{s.t.}\quad
% \sum_{i=1}^{K}w_{j,i}=K,\qquad
% 0\le w_{j,i}\le\frac{K}{k}.
% \]
% Let \(R_{j,(1)}\ge\cdots\ge R_{j,(K)}\),
% \(\tau_j=R_{j,(k)}\),
% \(n_{>,j}=|\{i:R_{j,i}>\tau_j\}|\), and
% \(n_{=,j}=|\{i:R_{j,i}=\tau_j\}|\). The weights are
% \[
% w_{j,i}^{\mathrm{top}\text{-}k}=\begin{cases}
% K/k, & R_{j,i}>\tau_j,\\[3pt]
% (K/k)(k-n_{>,j})/n_{=,j}, & R_{j,i}=\tau_j,\\[3pt]
% 0, & R_{j,i}<\tau_j.
% \end{cases}
% \]
% For \(K=24\), use \(k\in\{1,3,6,12\}\). With distinct rewards this is
% hard Top-$k$; only a tie at the boundary receives fractional weights. Both
% rules satisfy \(\sum_i w_{j,i}=K\). Compare performance across \(c\) and
% \(k\) to test how much concentration on the best samples helps.

\paragraph{Experiment B: Softmax with standard-deviation temperature.}

\[
f_c(R_{j,i})=\exp\!\left(\frac{R_{j,i}}{c\sigma_{c_j}}\right),
\qquad c\in\{0.5,0.75,1.0,1.5\}.
\]
For \(\sigma_{c_j}>0\), smaller \(c\) yields a steeper mapping and more
concentrated weights; larger \(c\) gives a flatter mapping.

\textbf{Target}: This exp tests
the trade-off between concentrating on high rewards and keeping a flatter
training signal. We want to test whether more concentrating or flat than classical softmax is better for optimize.

\paragraph{Experiment C: Linear baseline.}
\[
f_{\mathrm{lin}}(R)=[R]_+,
\qquad
f_{\mathrm{lin,shift}}(R)=[1+R]_+,
\qquad [x]_+=\max(x,0).
\]
Only \(f_{\mathrm{lin}}(R)=[R]_+\) is included in the planned sweep;
the shifted mapping is shown for comparison, not run as an experiment.

\textbf{Target}: Linear reward mapping acts as baseline variant. Expect to be worse than softmax and one-hot.

\paragraph{Experiment D: Sigmoid mapping.}
\[
f_{\mathrm{sig},c}(R_{j,i})
=\frac{1}{1+\exp(-(R_{j,i}-0.5)/c)},
\qquad c\in\{0.01,0.1,0.3\}.
\]

The smaller \(c\) is, the sharper the transition around \(R=0.5\).
After within-group normalization, the sigmoid may be sharper or flatter
than linear, depending on the reward range; its saturation also limits
how much the highest rewards can dominate.

\textbf{Target}: Compare sigmoid with the linear baseline within each group.
Test whether smaller \(c\) improves performance, whether the flatter
\(c=0.3\) setting underperforms, and whether a small-\(c\) sigmoid is a
better choice than the softmax-style mappings in Experiment B.


\paragraph{Experiment E: Power mapping.}

\[
f_p(R_{j,i})=\bigl[R_{j,i}\bigr]_+^p.
\]

We condider \(p\in\{0.5,2.0\}\).

\textbf{Target}: Compare with the linear \(p=1\) to test
whether the stronger emphasis on high rewards gives better behavior (similar
performance with softmax).

\paragraph{Experiment F: Softmax with L1 scale.}
Define
\[
||R_{j,i}-\bar{R}_{c_j}||_1=\frac{1}{K}\sum_{i=1}^{K}
\bigl|R_{j,i}-\bar R_{c_j}\bigr|.
\]
Then
\[
f_{c,1}(R_{j,i})
=\exp\!\left(\frac{R_{j,i}}{c||R_{j,i}-\bar{R}_{c_j}||_1}\right),
\qquad c\in\{0.5,0.75,1.0,1.5\}.
\]


\textbf{Target}: This is a variant for softmax. It is similiar with softmax with std temperature as shown in following figure. So we expect a similar performance with Experiment B. This exp can be skipped if nessesary.

\paragraph{Experiment G: Exponential on positive excess with mass-preserving rescaling.}
As in Experiment C, we use the linear mapping
\[
w_{j,i}=
\begin{cases}
\displaystyle
\frac{[R_{j,i}]_+}{\frac{1}{K}\sum_{k=1}^{K}[R_{j,k}]_+},
& \displaystyle \sum_{k=1}^{K}[R_{j,k}]_+>0,\\[10pt]
1, & \displaystyle \sum_{k=1}^{K}[R_{j,k}]_+=0,
\end{cases}
\]
so that \(\frac{1}{K}\sum_{i=1}^{K}w_{j,i}=1\). We
apply an exponential only to the positive excess, while leaving the
non-positive excess unchanged:
\[
h_{j,i}(\alpha)
=
\begin{cases}
\exp\!\bigl(\alpha\, (w_{j,i}-1)\bigr), & w_{j,i}-1>0,\\[4pt]
w_{j,i}-1, & w_{j,i}-1\le 0,
\end{cases}
\qquad
\alpha\in\{1.0,5.0,10.0,20.0\}.
\]
Because the exponential changes the total positive mass, we rescale the
positive part so that its sum is preserved:
\[
S_j^{+}=\sum_{i:\,w_{j,i}-1>0} (w_{j,i}-1),
\qquad
\widetilde{S}_j^{+}(\alpha)=\sum_{i:\,w_{j,i}-1>0} h_{j,i}(\alpha),
\]
\[
\tilde w_{j,i}
=
\begin{cases}
\displaystyle
\frac{S_j^{+}}{\widetilde{S}_j^{+}(\alpha)}\, h_{j,i}(\alpha),
& w_{j,i}-1>0,\\[10pt]
w_{j,i}-1, & w_{j,i}-1\le 0.
\end{cases}
\]
Since \(\sum_{i:\,w_{j,i}-1>0}\tilde w_{j,i}=S_j^{+}\) and the
negative part is untouched, we still have
\(\sum_{i=1}^{K}\tilde w_{j,i}=0\). The target shift then uses
the same form as the algorithm:
\[
\mathbf{x}_i^{\mathrm{target}}
=
\mathbf{x}_i^{\mathrm{old}}
+
\beta\,\tilde w_{j,i}(\mathbf{x}_i^0-\mathbf{x}_i^{\mathrm{old}}).
\]


\begin{figure}[thbp]
    \centering
    \includegraphics[width=1\linewidth]{imgs/weight_variants_normalized.png}
    \caption{Illustration of the temperature multiplier $c$ with
    $\bar R_{c_j}=0.3$, $\sigma_{c_j}=0.5$, and mean centered L1 deviation
    $d_{1,c_j}=0.4$.}
    \label{fig:placeholder}
\end{figure}

\begin{figure}[thbp]
    \centering
    \includegraphics[width=1\linewidth]{imgs/weight_excess_variants.png}
    \caption{Comparison Chart of Curves from Experiment G}
    \label{fig:placeholder}
\end{figure}




% \begin{figure}[H]
%     \centering
%     \includegraphics[width=1.0\linewidth]{imgs/softmax-variant-example.png}
%     \caption{Enter Caption}
%     \label{fig:placeholder}
% \end{figure}

\begin{table}[H]
    \centering
    \caption{Planned reward-to-weight experiments in execution order.}
    \label{tab:weight-variant-planned-sweep}
    \footnotesize
    \setlength{\tabcolsep}{5pt}
    \renewcommand{\arraystretch}{0.95}
    \begin{tabular}{@{}c l l c@{}}
        \toprule
        Run & Experiment & Configuration & Status \\
        \midrule
        01 & \multirow{4}{*}{A: near-max one-hot, $cR_j^{\max}$} & $c=0$  & - \\
        02 & & $c=0.05$ & - \\
        03 & & $c=0.1$ & - \\
        04 & & $c=0.25$ & - \\
        \midrule
        05 & \multirow{4}{*}{B: softmax, $c\sigma_{c_j}$} & $c=0.5$ & - \\
        06 & & $c=0.75$ & - \\
        07 & & $c=1.0$ & - \\
        08 & & $c=1.5$ & - \\
        \midrule
        09 & C: linear & $f(R)=[R]_+$ & - \\
        \midrule
        10 & \multirow{3}{*}{D: sigmoid, $(R-0.5)/c$} & $c=0.01$ & - \\
        11 & & $c=0.1$ & - \\
        12 & & $c=0.3$ & - \\
        \midrule
        13 & \multirow{2}{*}{E: power} & $f(R)=[R]_+^{0.5}$ & - \\
        14 & & $f(R)=[R]_+^{2.0}$ & - \\
        \midrule
        15 & \multirow{4}{*}{F: softmax, L1 scale} & $c=0.5$ & - \\
        16 & & $c=0.75$ & - \\
        17 & & $c=1.0$ & - \\
        18 & & $c=1.5$ & - \\
        \midrule
        19 & \multirow{4}{*}{G: positive excess} & $\alpha=1.0$ & - \\
        20 & & $\alpha=5.0$ & - \\
        21 & & $\alpha=10.0$ & - \\
        22 & & $\alpha=20.0$ & - \\
        \bottomrule
    \end{tabular}

    \vspace{0.3em}
    \scriptsize
    $\checkmark$ running\qquad - waiting\qquad $\bullet$ finished
\end{table}


\end{document}
